%% ============================================================
%%  SECCIÓN 7 — Resultados Experimentales
%% ============================================================

\section{Resultados Experimentales}
\label{sec:resultados}

En esta sección se presentan los datos crudos recolectados durante la sesión de laboratorio del 01 de octubre de 2026, así como los parámetros metrológicos constantes del montaje experimental (Cuadro~\ref{tab:parametros_constantes}).

\subsection{Parámetros Instrumentales Constantes}

\begin{table}[H]
    \centering
    \caption{Parámetros fijos del montaje experimental durante toda la sesión.}
    \label{tab:parametros_constantes}
    \begin{tabular}{@{}llc@{}}
        \toprule
        \textbf{Magnitud} & \textbf{Símbolo} & \textbf{Valor} \\
        \midrule
        Tensión de alimentación de la fuente DC & $V_s$ & \qty{7{,}50 \pm 0{,}05}{\volt} \\
        Incertidumbre del voltímetro digital & $\delta V$ & \qty{\pm 0{,}01}{\volt} \\
        Incertidumbre del amperímetro digital & $\delta I$ & \qty{\pm 0{,}001}{\ampere} \\
        Temperatura ambiente (estimada) & $T_{\text{amb}}$ & \qty{25 \pm 2}{\celsius} \\
        Valor nominal de la resistencia E2 & $R_{\text{nom}}$ & \qty{100}{\ohm} $\pm$ 5\,\% \\
        Límite de seguridad del diodo E3 & $I_{\text{max}}$ & \qty{0{,}900}{\ampere} \\
        \bottomrule
    \end{tabular}
\end{table}

\subsection{Tabla 1 --- Lámpara Incandescente de Tungsteno (Elemento E1)}

\begin{table}[H]
    \centering
    \caption{Mediciones directas $V$--$I$ de la lámpara incandescente (filamento de tungsteno). Se incluyen la resistencia estática $R = V/I$ y la potencia disipada $P = VI$ calculadas punto a punto.}
    \label{tab:datos_foquito}
    \resizebox{0.62\linewidth}{!}{%
    \begin{tabular}{ccccc}
        \toprule
        \textbf{Nro.} & \textbf{$V$ (\unit{\volt})} & \textbf{$I$ (\unit{\ampere})} & \textbf{$R = V/I$ (\unit{\ohm})} & \textbf{$P = VI$ (\unit{\watt})} \\
        \midrule
        1 & \num{0{,}00} & \num{0{,}000} & ---    & \num{0{,}000} \\
        2 & \num{1{,}00} & \num{0{,}148} & \num{6{,}76}  & \num{0{,}148} \\
        3 & \num{2{,}02} & \num{0{,}218} & \num{9{,}27}  & \num{0{,}440} \\
        4 & \num{3{,}04} & \num{0{,}272} & \num{11{,}18} & \num{0{,}827} \\
        5 & \num{4{,}01} & \num{0{,}318} & \num{12{,}61} & \num{1{,}275} \\
        6 & \num{5{,}01} & \num{0{,}366} & \num{13{,}69} & \num{1{,}834} \\
        7 & \num{6{,}01} & \num{0{,}412} & \num{14{,}59} & \num{2{,}476} \\
        8 & \num{6{,}43} & \num{0{,}432} & \num{14{,}88} & \num{2{,}778} \\
        \bottomrule
    \end{tabular}%
    }
\end{table}

\subsection{Tabla 2 --- Resistencia de Composición de Carbón (Elemento E2)}

\begin{table}[H]
    \centering
    \caption{Mediciones directas $V$--$I$ de la resistencia de composición de carbón de \qty{100}{\ohm}. Se incluyen $R = V/I$ y su incertidumbre propagada $\delta R$ (Ec.~\eqref{eq:prop_errores_R}).}
    \label{tab:datos_resistencia}
    \resizebox{0.75\linewidth}{!}{%
    \begin{tabular}{cccccc}
        \toprule
        \textbf{Nro.} & \textbf{$V$ (\unit{\volt})} & \textbf{$I$ (\unit{\ampere})} & \textbf{$R = V/I$ (\unit{\ohm})} & \textbf{$\delta R$ (\unit{\ohm})} & \textbf{$E_r$ (\unit{\percent})} \\
        \midrule
        1 & \num{0{,}00} & \num{0{,}000} & --- & ---   & ---  \\
        2 & \num{1{,}00} & \num{0{,}010} & \num{100{,}0} & \num{\pm 11{,}0} & \num{0{,}0} \\
        3 & \num{2{,}00} & \num{0{,}020} & \num{100{,}0} & \num{\pm 5{,}5}  & \num{0{,}0} \\
        4 & \num{3{,}05} & \num{0{,}031} & \num{98{,}4}  & \num{\pm 3{,}5}  & \num{1{,}6} \\
        5 & \num{4{,}03} & \num{0{,}040} & \num{100{,}8} & \num{\pm 2{,}8}  & \num{0{,}8} \\
        6 & \num{5{,}00} & \num{0{,}050} & \num{100{,}0} & \num{\pm 2{,}2}  & \num{0{,}0} \\
        7 & \num{6{,}02} & \num{0{,}060} & \num{100{,}3} & \num{\pm 1{,}8}  & \num{0{,}3} \\
        8 & \num{6{,}49} & \num{0{,}065} & \num{99{,}8}  & \num{\pm 1{,}7}  & \num{0{,}2} \\
        \bottomrule
    \end{tabular}%
    }
\end{table}

Todos los valores individuales de $R = V/I$ son compatibles con el valor nominal de \qty{100}{\ohm} dentro de su incertidumbre instrumental $\delta R$.

\subsection{Tabla 3 --- Diodo Semiconductor de Silicio (Elemento E3)}

\begin{table}[H]
    \centering
    \caption{Mediciones directas $V$--$I$ del diodo de silicio en polarización directa. Se incluyen $R = V/I$ y $P = VI$. El punto \texorpdfstring{N.\textsuperscript{o}\,4}{Nro. 4} (\qty{0{,}70}{\volt}; \qty{0{,}230}{\ampere}) se identifica como punto atípico (*) por el análisis residual.}
    \label{tab:datos_diodo}
    \resizebox{0.62\linewidth}{!}{%
    \begin{tabular}{ccccc}
        \toprule
        \textbf{Nro.} & \textbf{$V$ (\unit{\volt})} & \textbf{$I$ (\unit{\ampere})} & \textbf{$R = V/I$ (\unit{\ohm})} & \textbf{$P = VI$ (\unit{\watt})} \\
        \midrule
        1 & \num{0{,}00} & \num{0{,}000} & ---          & \num{0{,}000} \\
        2 & \num{0{,}61} & \num{0{,}050} & \num{12{,}2} & \num{0{,}031} \\
        3 & \num{0{,}69} & \num{0{,}215} & \num{3{,}21} & \num{0{,}148} \\
        4* & \num{0{,}70} & \num{0{,}230} & \num{3{,}04} & \num{0{,}161} \\
        5 & \num{0{,}71} & \num{0{,}411} & \num{1{,}73} & \num{0{,}292} \\
        6 & \num{0{,}72} & \num{0{,}537} & \num{1{,}34} & \num{0{,}387} \\
        7 & \num{0{,}73} & \num{0{,}690} & \num{1{,}06} & \num{0{,}504} \\
        8 & \num{0{,}75} & \num{0{,}783} & \num{0{,}96} & \num{0{,}587} \\
        \bottomrule
    \end{tabular}%
    }
\end{table}

(*) Punto identificado como atípico mediante el análisis de residuos respecto al modelo exponencial de Shockley. Se reporta tal como se midió; su diagnóstico cuantitativo se desarrolla en la Sección~\ref{sec:analisis_exp} y en la Sección~\ref{sec:tratamiento}.